%!TeX root = ../main.tex

% ==================================================================================================================================
% 6. Metodologia
% ==================================================================================================================================
\section{Metodologia}

Este capítulo descreve o ambiente de testes construído para \textit{implementar}
e \textit{comparar} os algoritmos dos capítulos anteriores. São apresentados a
arquitetura do ambiente, a interface padrão de cada problema, a origem e a
adaptação de cada algoritmo, os geradores de instância, o protocolo de medição,
os critérios de comparação e a máquina utilizada, de modo que os resultados do
capítulo seguinte sejam reprodutíveis e as comparações, defensáveis.

% ---------------------------------------------------------------------------
\subsection{Arquitetura do ambiente de testes}

O ambiente é organizado em três camadas, com separação estrita entre código de
terceiros e código próprio. Na primeira, cada algoritmo obtido de um repositório
público é mantido \textbf{intacto}, garantindo que o código medido é exatamente
o publicado por seus autores --- condição essencial para a validade das
comparações e para a análise de eventuais defeitos. Na segunda, um
\textbf{adaptador} (\textit{wrapper}) por algoritmo traduz sua interface
particular para a \textbf{interface padrão} do problema, documentando em seu
cabeçalho a complexidade, o tipo do método e a fonte de origem. Na terceira, um
único \textbf{módulo de medição} --- responsável por geração de instâncias,
cronometragem, verificação e gravação --- é compartilhado por todos os problemas,
o que garante que todos os algoritmos sejam submetidos ao mesmo protocolo. A
separação entre as duas primeiras camadas permite ainda comparar, sob condições
idênticas, múltiplas implementações de um mesmo algoritmo obtidas de autores
diferentes.

% ---------------------------------------------------------------------------
\subsection{Interface padrão dos problemas}

Cada problema possui uma interface padrão à qual todos os seus algoritmos são
adaptados (Tabela~\ref{tab:interfaces}). É ela que torna a comparação justa: todo
algoritmo de um mesmo problema recebe exatamente a mesma instância e produz a
saída no mesmo formato, de modo que a única variável entre eles é o próprio
método.

\begin{table}[H]
    \centering
    \caption{Interface padrão de entrada e saída por problema}
    \label{tab:interfaces}
    \renewcommand{\arraystretch}{1.3}
    \footnotesize
    \begin{tabularx}{\textwidth}{@{} l X X @{}}
        \toprule
        \textbf{Problema} & \textbf{Entrada} & \textbf{Saída} \\
        \midrule
        Cardinalidade máxima (bipartido) & vértices de cada lado e lista de arestas & tamanho e emparelhamento \\
        Atribuição linear & matriz de custos $n \times n$ & custo total e atribuição \\
        Atribuição gargalo & matriz de custos $n \times n$ & custo gargalo e atribuição \\
        Casamento estável & listas de preferência de ambos os lados & emparelhamento \\
        Cardinalidade máxima (geral) & número de vértices e lista de arestas & tamanho e emparelhamento \\
        Peso máximo (geral) & número de vértices e arestas com peso & peso total e emparelhamento \\
        Emparelhamento induzido máximo & número de vértices e lista de arestas & tamanho e emparelhamento \\
        \bottomrule
    \end{tabularx}
\end{table}

% ---------------------------------------------------------------------------
\subsection{Geradores de instância}

As instâncias são geradas por construtores parametrizados pelo tamanho $n$ e por
uma semente fixa, o que as torna determinísticas e reprodutíveis. Para cada
problema definiram-se famílias com características distintas, de modo a expor o
comportamento dos algoritmos sob diferentes estruturas de entrada: para
cardinalidade máxima em grafos bipartidos, grafos esparsos (com e sem
emparelhamento perfeito) e densos; para atribuição linear e gargalo, matrizes
densas aleatórias, matrizes estruturadas (custo em função da distância entre
índices) e matrizes com poucos valores distintos (muitos empates); para
casamento estável, preferências aleatórias, correlacionadas (\textit{ranking} de
popularidade compartilhado, que gera contenção) e cíclicas; e, para grafos gerais
e emparelhamento induzido máximo, grafos aleatórios de densidade controlada,
ponderados quando aplicável.

% ---------------------------------------------------------------------------
\subsection{Protocolo de medição}
\label{sec:protocolo}

A medição de tempo é sensível a fatores externos ao próprio método; o protocolo a
seguir reduz essas fontes de ruído. Uma \textbf{execução de aquecimento} é
descartada, por pagar custos não recorrentes (importações, alocação inicial,
caches frios); ela serve apenas para aferir o orçamento de tempo e fornecer o
resultado usado na verificação de qualidade. O \textbf{coletor de lixo é
desligado} durante as repetições cronometradas --- para que uma coleta não
apareça como lentidão do algoritmo --- e reativado entre as células. Cada célula
é medida em várias repetições, das quais se reporta a \textbf{mediana}, robusta a
interrupções momentâneas do sistema, ao contrário da média. Todos os algoritmos
de um problema recebem \textbf{as mesmas instâncias}. Por fim, um \textbf{limite
de tempo por célula} impede que um único algoritmo defeituoso inviabilize toda a
suíte --- situação efetivamente observada e discutida no capítulo de resultados:
se o aquecimento excede o limite, a célula é marcada como estourada e os tamanhos
maiores daquele algoritmo são omitidos.

Registra-se uma limitação deliberada: o limite de tempo é aplicado em torno da
execução de aquecimento, não a cada operação. Uma célula patológica é
interrompida ao ultrapassar o limite, mas não instantaneamente --- escolha que
evita o sobrecusto de uma verificação passo a passo na própria região
cronometrada.

% ---------------------------------------------------------------------------
\subsection{Critérios de comparação}
\label{sec:criterios}

Comparar tempo de execução entre implementações em linguagens diferentes é
intrinsecamente injusto: parte da diferença reflete a linguagem, não o algoritmo.
Para que as comparações sejam válidas, adotaram-se quatro critérios
complementares.

Primeiro, as implementações em linguagem interpretada e as de núcleo compilado
são apresentadas em \textbf{tabelas separadas}, nunca num mesmo ranking de tempo;
as de núcleo compilado servem para validar a corretude e indicar a ordem de
grandeza do estado da prática. Vale notar que nem toda biblioteca é compilada:
algumas bibliotecas maduras usadas como referência são escritas na mesma
linguagem interpretada dos demais algoritmos, e nesses casos a comparação de
tempo contra elas é legítima. Há ainda uma zona intermediária --- algoritmos que
são laços na linguagem interpretada operando sobre vetores de uma biblioteca
numérica de núcleo compilado; classificam-se como interpretados, com essa
característica declarada. Segundo, as tabelas trazem a \textbf{complexidade
assintótica} de tempo e espaço ao lado do tempo medido, permitindo confrontar o
previsto com o observado. Terceiro, para cada algoritmo estima-se o
\textbf{expoente empírico de crescimento} --- a inclinação da reta ajustada ao
gráfico do logaritmo do tempo pelo logaritmo do tamanho ---, que, ao contrário do
tempo absoluto, é \textbf{comparável entre linguagens}, pois a diferença de
linguagem afeta a constante (a posição da reta), não a inclinação; reporta-se
também a qualidade do ajuste.

Nos problemas com apenas algoritmos exatos, a qualidade da solução é idêntica
entre eles e a comparação recai sobre o tempo. Nos problemas com heurísticas,
cada heurística reporta a \textbf{razão} entre a qualidade que encontra e a do
algoritmo exato de referência, caracterizando o compromisso entre qualidade e
tempo.

% ---------------------------------------------------------------------------
\subsection{Ambiente de execução e reprodutibilidade}
\label{sec:ambiente}

Todos os resultados foram obtidos numa única máquina
(Tabela~\ref{tab:maquina}); medir tudo no mesmo ambiente é o que torna os tempos
comparáveis. Três características dela exigiram cuidados, todos documentados para
permitir reprodução. Por ser um \textbf{processador híbrido} (núcleos de
desempenho e de eficiência, não equivalentes), a execução foi fixada a um núcleo
de desempenho, evitando que a migração entre núcleos alterasse o tempo sem
mudança no código. Como a \textbf{frequência varia} dinamicamente com a carga, o
regulador foi fixado no modo de desempenho, com a máquina conectada à energia e
sem aplicações concorrentes. E, como a \textbf{ordem de iteração de coleções de
objetos} em Python é aleatorizada por processo --- e um dos algoritmos se mostrou
sensível a ela, conforme o capítulo de resultados ---, a semente de aleatorização
foi fixada, tornando as execuções reprodutíveis.

\begin{table}[H]
    \centering
    \caption{Especificação da máquina de medição}
    \label{tab:maquina}
    \renewcommand{\arraystretch}{1.3}
    \footnotesize
    \begin{tabularx}{\textwidth}{@{} l X @{}}
        \toprule
        \textbf{Componente} & \textbf{Especificação} \\
        \midrule
        Processador & Intel Core 7 150U (10 núcleos físicos, 12 \textit{threads}) \\
        Topologia & híbrida: 2 núcleos de desempenho e 8 de eficiência \\
        Memória & 33 GB \\
        Sistema operacional & Ubuntu 24.04.2 LTS (kernel 7.0.0) \\
        Linguagem & Python 3.12.3 \\
        Bibliotecas & NumPy 2.5.2, SciPy 1.18.0, NetworkX 3.6.1, igraph 1.0.0, lap 0.5.13 \\
        \bottomrule
    \end{tabularx}
\end{table}

% ---------------------------------------------------------------------------
\subsection{Origem, adaptação e verificação dos algoritmos}
\label{sec:proveniencia}

Em coerência com o objetivo de comparar métodos \textit{presentes na literatura},
a maior parte dos algoritmos foi obtida de implementações públicas, e não
reescrita: adaptar código consolidado reduz o risco de erros próprios e permite
avaliar a qualidade das implementações efetivamente disponíveis. Conforme a
natureza do código de origem, empregaram-se quatro estratégias de adaptação:
\textbf{(i)} carregamento direto, quando o original é utilizável na linguagem do
projeto; \textbf{(ii)} encapsulamento, para métodos disponíveis em bibliotecas
maduras; \textbf{(iii)} reimplementação fiel, quando o original não é utilizável
diretamente --- por estar em outra linguagem ou versão, ou por conter efeitos
colaterais ---, seguindo a lógica original com a fonte documentada; e
\textbf{(iv)} implementação própria a partir da fonte bibliográfica, para as
heurísticas e alguns exatos sem implementação pública adequada, algoritmos de
descrição curta cuja reimplementação é mais confiável que um repositório de
qualidade incerta. A Tabela~\ref{tab:proveniencia} registra o tipo e a origem de
cada algoritmo.

\begin{table}[H]
    \centering
    \caption{Proveniência dos algoritmos avaliados}
    \label{tab:proveniencia}
    \renewcommand{\arraystretch}{1.2}
    \footnotesize
    \begin{tabularx}{\textwidth}{@{} l l X @{}}
        \toprule
        \textbf{Algoritmo} & \textbf{Tipo} & \textbf{Origem} \\
        \midrule
        \multicolumn{3}{@{}l}{\textit{Cardinalidade máxima --- grafos bipartidos}} \\
        Caminhos aumentantes & Exato & repositório \texttt{wbchristerson/perfect-matchings} \\
        Hopcroft-Karp & Exato & repositório \texttt{sofiat-olaosebikan/hopcroftkarp} \\
        Edmonds-Karp (fluxo) & Exato & repositório \texttt{Maxflow-Algorithms} (reimplementado) \\
        Hopcroft-Karp (NetworkX) & Referência & biblioteca NetworkX (interpretada) \\
        Emparelhamento bipartido (igraph) & Referência & biblioteca igraph (núcleo compilado) \\
        Gulosa & Heurística & implementação própria \cite{cormen} \\
        Karp-Sipser & Heurística & implementação própria \cite{karp-sipser} \\
        \addlinespace
        \multicolumn{3}{@{}l}{\textit{Atribuição linear}} \\
        Húngaro (munkres) & Exato & repositório \texttt{bmc/munkres} \\
        Húngaro (benchaplin) & Exato & repositório \texttt{benchaplin/hungarian-algorithm} \\
        Húngaro (hunalgorithm) & Exato & repositório \texttt{pbertoni/HungarianAlgorithm} \\
        Húngaro (mayorx) & Exato & repositório \texttt{mayorx/hungarian-algorithm} \\
        Caminhos mínimos sucessivos & Exato & reimplementado a partir de repositório \\
        SciPy & Referência & biblioteca SciPy (núcleo compilado) \\
        Jonker-Volgenant & Referência & biblioteca lap (núcleo compilado) \\
        \addlinespace
        \multicolumn{3}{@{}l}{\textit{Atribuição gargalo}} \\
        Threshold & Exato & reimplementado (original em JavaScript) \\
        Caminhos aumentantes & Exato & implementação própria \cite{burkard} \\
        \addlinespace
        \multicolumn{3}{@{}l}{\textit{Casamento estável}} \\
        Gale-Shapley (orientado a objetos) & Exato & pacote \texttt{gale-shapley-algorithm} \\
        Gale-Shapley (numérico) & Exato & pacote \texttt{gale-shapley-algorithm} \\
        Gale-Shapley (Lattas) & Exato & reimplementado (repositório \texttt{StableMarriages}) \\
        \addlinespace
        \multicolumn{3}{@{}l}{\textit{Grafos gerais}} \\
        Blossom/Edmonds (cardinalidade) & Exato & repositório \texttt{adharshkamath/Edmonds-Algorithm} \\
        van Rantwijk (peso máximo) & Exato & \texttt{mwmatching.py} de J.\ van Rantwijk \\
        NetworkX (cardinalidade e peso) & Referência & biblioteca NetworkX (interpretada) \\
        Path Growing & Heurística & implementação própria \cite{drake} \\
        \addlinespace
        \multicolumn{3}{@{}l}{\textit{Emparelhamento induzido máximo}} \\
        Exato (conjunto independente) & Exato & implementação própria \cite{cameron} \\
        Guloso aleatório & Heurística & implementação própria \\
        Guloso de grau mínimo & Heurística & implementação própria \\
        \bottomrule
    \end{tabularx}
\end{table}

Os endereços completos constam no cabeçalho de cada adaptador, no código-fonte.
Cabe uma ressalva de linhagem: a rotina de emparelhamento de peso máximo em
grafos gerais do NetworkX é uma adaptação do mesmo código de van Rantwijk também
avaliado de forma independente; compará-las mede o efeito da engenharia de
biblioteca sobre um mesmo algoritmo, e não duas abordagens distintas.

Por fim, nenhuma implementação própria ou reimplementada foi incorporada sem
\textbf{verificação independente}: os métodos exatos foram confrontados com
soluções por força bruta (em instâncias pequenas) e com bibliotecas
consolidadas; as heurísticas, quanto à validade da solução, a não superarem o
ótimo e a respeitarem suas garantias de aproximação. Além disso, os benchmarks
já existentes ao início do trabalho foram migrados para o ambiente unificado sob
uma salvaguarda de não-regressão: a qualidade de saída de cada um foi registrada
antes da migração e exigiu-se que a nova execução a reproduzisse exatamente, de
modo que a unificação não alterasse nenhuma conclusão.
